\section{Conclusion and Discussion}
\label{sec:conclusion}

In this work, we have executed a rigorous empirical deconstruction of Quantum Wavefunction Superposition Merging (QWS-Merge), a recent state-of-the-art quantum-inspired model-merging framework. From our perspective as \textbf{The Methodologist}, we identified key confounding variables in its original evaluation—namely, the omission of basic classical baselines and the lack of a controlled representation evaluation. By isolating the structural components of routing under identical low-dimensional projections, we proposed a simple classical alternative: the Layer-wise Low-dimensional Classical Router (L3-Router).

Our experimental results within a controlled representation sandbox demonstrate a critical methodological truth. The wave-inspired SOTA QWS-Merge completely collapses, achieving a Joint Mean of only \textbf{36.10\%}—performing worse than static uniform merging (\textbf{43.40\%}). By replacing the wave formulation with our proposed classical \textbf{L3-Linear} router, we avoid this collapse and achieve a Joint Mean of \textbf{63.10\%} (\textbf{+27.00\%} over QWS-Merge). Most remarkably, our deconstruction reveals that the simplest baseline of all—the global, unregularized classical \textbf{Linear Router}—outperforms all other models, achieving the highest Joint Mean of \textbf{67.20\%}. This serves as a profound methodological warning for the deep learning community: before introducing complex, multi-layer specialized routing or wrapping models in quantum analogies, researchers must first rigorously evaluate simple global classical baselines.

Furthermore, our deployment stream audit reveals that dynamic model merging suffers from a severe \textit{heterogeneity collapse} when subjected to mixed-task inference batches, dropping the Linear Router to \textbf{51.10\%} and QWS-Merge to \textbf{10.80\%}. We critically deconstruct our own proposed \textbf{L3-Softmax} router under these mixed-task streams, exposing the \textbf{``Robustness-Accuracy Illusion''}. Although L3-Softmax exhibits relative stability under stream shifts (dropping by only \textbf{4.10\%} compared to the Linear Router's \textbf{16.10\%} drop), its absolute accuracy is inferior in all deployment scenarios. This reveals that its apparent robustness is an artifact of the Softmax simplex constraint forcing dynamic routing coefficients toward a mediocre, uniform-like average, serving as a cautionary tale against celebrating relative stability metrics over absolute baseline performance.

Crucially, our closed-form proof of layer-averaging collapse (Section~3.5) is not restricted to our specific sandbox setup. Rather, it applies universally to \textit{any} dynamic routing model that attempts to merge a shared, global head or global parameter group (such as a shared classification head, an embedding layer, or a final language model projection head) using layer-wise specialized routing coefficients. Because averaging layer-wise routing weights collapses the multi-layer routing parameter space back to a single-layer routing space, any architectural specialized complexity across layers becomes mathematically redundant and introduces optimization noise. This proof serves as a broad, foundational warning for future architecture design in parameter merging, urging researchers to mathematically analyze potential parameter collapses before claiming multi-layer specialized capacity.

To bridge the gap between our controlled sandbox and practical production-scale systems, we outline a concrete, immediately actionable deployment roadmap on Vision-Language CLIP-ViT-B and Large Language Models (LLMs) in Appendix~\ref{sec:roadmap_app}. We provide step-by-step mathematical algorithms for unsupervised PCA and zero-shot text prompt projections. We also detail the hardware-grounded trade-offs of compiler-level parallel execution (including custom Triton-based dynamic weight assembly kernels, which incur exactly $2 K \cdot M$ multiply-accumulate FLOPs per layer and scale memory bandwidth by $(1 + K \cdot \gamma) M$ bytes using LoRA parameterization) to enable practitioners to perform precise, hardware-informed cost-benefit analyses before real-world dynamic deployment. Specifically, when scaling to massive generative autoregressive LLMs (such as LLaMA-3-8B or Mistral-7B), the L3-Router can be deployed seamlessly alongside quantized low-rank task adapters (LoRAs) rather than full-scale backbone parameter matrices. By parameterizing downstream tasks as low-rank offset updates and utilizing custom Triton-level parallel kernels to dynamically assemble the active routing path on-the-fly, practitioners can completely bypass physical storage limits and deploy dozens of specialized LLM experts simultaneously on a single consumer GPU. We also detail our study's limitations and discuss projection dimension hyperparameter sensitivity in Appendix~\ref{sec:limitations_app}.

In conclusion, our findings serve as a strong cautionary tale for the deep learning community. True scientific progress is not driven by wrapping standard neural layers in complex mathematical analogies or quantum mechanics vocabulary. Instead, it is achieved through systematic deconstruction, rigorous baseline tuning, and a persistent commitment to experimental hygiene. We encourage future model-merging research to prioritize simple, transparent, and properly regularized baselines before claiming superiority through over-engineered architectural complexity.
